% Figure from MATH 590 notes, section 24. Compile with the notes preamble.
\begin{tikzpicture}[
  hx/.style={variable=\t, samples=30, smooth},
  op/.style={circle, inner sep=1.3pt, fill=white, draw=figR, thick}]
  \def\R{1.5}\def\r{0.2}\def\h{1.25}\def\yo{1.3}
  % screen position of the helix point over parameter x: (R sin 2pi x, -r cos 2pi x + h x + yo)
  % back halves (n+1/4, n+3/4): faint blue
  \foreach \n in {-2,-1,0,1,2} {
    \draw[figB!60, thick, hx, domain={max(\n+0.25,-1.4)}:{min(\n+0.75,2.4)}]
      plot ({\R*sin(360*\t)}, {-\r*cos(360*\t)+\h*\t+\yo});
  }
  % front halves = slices V_n = (n-1/4, n+1/4): red
  \foreach \n in {-1,0,1,2} {
    \draw[figR, very thick, hx, domain=\n-0.25:\n+0.25]
      plot ({\R*sin(360*\t)}, {-\r*cos(360*\t)+\h*\t+\yo});
    \node[op] at ({-\R}, {\h*(\n-0.25)+\yo}) {};
    \node[op] at ({\R}, {\h*(\n+0.25)+\yo}) {};
  }
  % fiber over b0
  \draw[black!40, dotted, thin] (0,-1.55) -- (0,{\h*2-\r+\yo});
  \foreach \n in {-1,0,1,2} {
    \fill (0,{\h*\n-\r+\yo}) circle (1.5pt);
    \node[font=\scriptsize, anchor=south west, inner sep=1.5pt] at (0.02,{\h*\n-\r+\yo+0.02}) {$\n$};
  }
  \foreach \n in {0,1} {
    \node[font=\scriptsize, figR, anchor=west] at ({\R+0.1}, {\h*(\n+0.25)+\yo}) {$V_{\n}$};
  }
  \node[font=\small, figB] at (-2.1,{\h*2+\yo}) {$\mathbb{R}$};
  \draw[-stealth, thick, black!70] (-2.1,-0.35) -- node[left, font=\scriptsize] {$p$} (-2.1,-1.1);
  % base circle
  \begin{scope}[yshift=-1.55cm]
    \draw[figB!60, thick] (\R,0) arc[start angle=0, end angle=180, x radius=\R, y radius=\r];
    \draw[figR, very thick] (-\R,0) arc[start angle=180, end angle=360, x radius=\R, y radius=\r];
    \node[op] at (-\R,0) {};
    \node[op] at (\R,0) {};
    \fill (0,-\r) circle (1.5pt);
    \node[font=\scriptsize, below, inner sep=2pt] at (0,-\r) {$b_0 = (1,0)$};
    \node[font=\scriptsize, figR, anchor=west] at (\R+0.1,0) {$U$};
    \node[font=\small, figB] at (-2.1,0) {$S^1$};
  \end{scope}
\end{tikzpicture}
